\chapter{\'algebra Relacional de JOINs, Subconsultas y CTEs}
\label{cap:algebra_joins_subconsultas}

\section{Marco Conceptual y Te\'orico (Curr\'iculo Universitario)}

El \'algebra Relacional provee la base matem\'atica de las operaciones de combinaci\'on de relaciones.

\subsection{Familias de Operadores JOIN}
\begin{itemize}
    \item \textbf{Inner Join ($R \bowtie S$):} Intersecci\'on de tuplas que satisfacen la condici\'on de correspondencia.
    \item \textbf{Left Outer Join ($R \bowtie_{L} S$):} Preserva todas las tuplas de la relaci\'on izquierda $R$, rellenando con \texttt{NULL} cuando no hay coincidencia en $S$.
    \item \textbf{Right Outer Join ($R \bowtie_{R} S$):} Preserva todas las tuplas de la relaci\'on derecha $S$.
    \item \textbf{Full Outer Join ($R \bowtie_{F} S$):} Preserva todas las tuplas de ambas relaciones.
    \item \textbf{Cross Join ($R \times S$):} Producto cartesiano exhaustivo de todas las tuplas.
\end{itemize}

---

\section{Taller de Consolidaci\'on del Cap\'itulo}

\begin{businessproblem}
El Director de Ventas necesita un reporte que consolide el total facturado por cliente durante el a\~no 1997, incluyendo el nombre de la compa\~n\'ia, pa\'is y el n\'umero total de \'ordenes realizadas, ordenado de mayor a menor facturaci\'on y asegurando que clientes sin compras aparezcan con monto cero.
\end{businessproblem}

\subsection{Soluci\'on T\'ecnica en PostgreSQL con CTEs y LEFT JOIN}

\begin{lstlisting}[style=sqlstyle, caption={Consulta modular con CTE y agregaci\'on sobre Northwind}]
WITH customer_order_totals AS (
    SELECT
        o.customer_id,
        COUNT(DISTINCT o.order_id) AS total_orders,
        SUM(od.unit_price * od.quantity * (1 - od.discount)) AS total_spent
    FROM orders AS o
    INNER JOIN order_details AS od
        ON o.order_id = od.order_id
    WHERE EXTRACT(YEAR FROM o.order_date) = 1997
    GROUP BY
        o.customer_id
)
SELECT
    c.customer_id,
    c.company_name,
    c.country,
    COALESCE(cot.total_orders, 0) AS total_orders_1997,
    COALESCE(cot.total_spent, 0.00)::NUMERIC(12, 2) AS total_revenue_1997
FROM customers AS c
LEFT JOIN customer_order_totals AS cot
    ON c.customer_id = cot.customer_id
ORDER BY
    total_revenue_1997 DESC,
    c.company_name ASC;
\end{lstlisting}

\begin{engtip}
\begin{itemize}
    \item \textbf{Estructura con CTE (\texttt{WITH}):} Desacopla la l\'ogica de agregaci\'on del detalle de la uni\'on con la tabla de clientes, mejorando la legibilidad y permitiendo al optimizador ejecutar agregaciones tempranas antes del join final.
\end{itemize}
\end{engtip}
